
\documentclass[11pt,a4paper]{article}

\usepackage[margin=1in]{geometry}
\usepackage{amsmath,amssymb}
\usepackage{siunitx}
\usepackage{booktabs}
\usepackage{array}
\usepackage{enumitem}
\usepackage{hyperref}
\usepackage{fancyhdr}

\hypersetup{
    colorlinks=true,
    linkcolor=blue,
    citecolor=blue,
    urlcolor=blue
}

\pagestyle{fancy}
\setlength{\headheight}{14pt}
\fancyhf{}
\lhead{Quantum Mechanics 3}
\rhead{Quantum Mechanics 3}
\cfoot{\thepage}

\title{\textbf{Quantum Mechanics 3}\\Worked Solutions}
\author{Worked Solutions}
\date{}

\begin{document}

\maketitle

\section*{Introduction}

This document presents worked solutions to the \emph{Quantum Mechanics 3} problem
sheet supplied with the BPhO Computational Physics Challenge 2026. The sheet
focuses on de Broglie matter waves, the uncertainty principle, wavefunctions,
the particle in a box, and the Copenhagen interpretation of measurement.

The derivations and numerical methods follow the accompanying BPhO solution
document. Numerical values are quoted to a sensible precision and use the
physical constants supplied on the problem sheet.

% ============================================================
\section{Question 1 --- Matter waves, uncertainty and measurement}

\subsection{(i) de Broglie wavelength of an electron}

For a non-relativistic particle,
\[
p=m_ev,
\]
and the de Broglie relation is
\[
\lambda=\frac{h}{p}.
\]

The electron travels at
\[
v=\frac{c}{137}.
\]

Therefore
\[
\lambda
=
\frac{h}{m_ec/137}
=
\frac{137h}{m_ec}.
\]

Using
\[
h=6.626\times10^{-34}\ \mathrm{J\,s},
\qquad
m_e=9.109\times10^{-31}\ \mathrm{kg},
\qquad
c=2.998\times10^8\ \mathrm{m\,s^{-1}},
\]
gives
\[
\lambda
=
\frac{137(6.626\times10^{-34})}
{(9.109\times10^{-31})(2.998\times10^8)}
\approx3.3\times10^{-10}\ \mathrm m.
\]

Thus
\[
\boxed{\lambda\approx3.3\times10^{-10}\ \mathrm m}.
\]

The supplied solution also checks that using the relativistic momentum makes no
significant difference at this speed to the quoted precision.

\subsection{(ii) Equal de Broglie wavelengths for an electron and neutron}

If
\[
\lambda_e=\lambda_n,
\]
then
\[
\frac{h}{p_e}=\frac{h}{p_n},
\]
so
\[
p_e=p_n.
\]

Using classical momentum,
\[
m_ev_e=m_nv_n.
\]

Therefore
\[
\boxed{
\frac{v_e}{v_n}=\frac{m_n}{m_e}
}.
\]

Numerically,
\[
\frac{v_e}{v_n}
=
\frac{1.6749\times10^{-27}}
{9.109\times10^{-31}}
\approx1837.
\]

Hence
\[
\boxed{\frac{v_e}{v_n}\approx1837}.
\]

The classical kinetic energy is
\[
E=\frac{p^2}{2m}.
\]

Since the momenta are equal,
\[
\frac{E_e}{E_n}
=
\frac{m_n}{m_e}
\approx1837.
\]

Thus
\[
\boxed{\frac{E_e}{E_n}\approx1837}.
\]

So for equal de Broglie wavelength, the electron travels about $1837$ times
faster and has about $1837$ times the kinetic energy of the neutron.

\subsection{(iii) Range and lifetime of a $Z$ boson}

The problem gives
\[
\Delta x\,\Delta p\approx\frac{\hbar}{2},
\qquad
\Delta E\,\Delta t\approx\frac{\hbar}{2},
\]
with
\[
\Delta p\approx mc,
\qquad
\Delta E\approx mc^2.
\]

Hence
\[
\Delta x\approx\frac{\hbar}{2mc},
\]
and
\[
\Delta t\approx\frac{\hbar}{2mc^2}.
\]

The mass of the $Z$ boson is
\[
m=\frac{91.2\ \mathrm{GeV}}{c^2}.
\]

Converting this to kilograms,
\[
m
=
\frac{91.2\times10^9(1.602\times10^{-19})}
{(2.998\times10^8)^2}
\approx1.626\times10^{-25}\ \mathrm{kg}.
\]

Therefore
\[
\Delta x
\approx
\frac{6.626\times10^{-34}}
{2(1.626\times10^{-25})(2.998\times10^8)}
\approx1.08\times10^{-17}\ \mathrm m.
\]

Thus
\[
\boxed{\Delta x\approx1.08\times10^{-17}\ \mathrm m}.
\]

Similarly,
\[
\Delta t
\approx
\frac{6.626\times10^{-34}}
{2(1.626\times10^{-25})(2.998\times10^8)^2}
\approx3.61\times10^{-27}\ \mathrm s.
\]

Hence
\[
\boxed{\Delta t\approx3.61\times10^{-27}\ \mathrm s}.
\]

These scales are far below ordinary atomic dimensions and timescales.

\subsection{(iv) Verifying a solution of the wave equation}

Consider
\[
\psi(x,t)=A\sin(kx-\omega t).
\]

The spatial derivatives are
\[
\frac{\partial\psi}{\partial x}
=
kA\cos(kx-\omega t),
\]
and therefore
\[
\frac{\partial^2\psi}{\partial x^2}
=
-k^2A\sin(kx-\omega t)
=
-k^2\psi.
\]

The time derivatives are
\[
\frac{\partial\psi}{\partial t}
=
-\omega A\cos(kx-\omega t),
\]
and
\[
\frac{\partial^2\psi}{\partial t^2}
=
-\omega^2A\sin(kx-\omega t)
=
-\omega^2\psi.
\]

Now
\[
k=\frac{2\pi}{\lambda},
\qquad
\omega=2\pi f,
\qquad
c=f\lambda.
\]

Thus
\[
\omega=ck,
\]
so
\[
\omega^2=c^2k^2.
\]

Therefore
\[
\frac{\partial^2\psi}{\partial t^2}
=
-c^2k^2\psi
=
c^2\frac{\partial^2\psi}{\partial x^2}.
\]

Hence
\[
\boxed{
\frac{\partial^2\psi}{\partial x^2}
=
\frac{1}{c^2}
\frac{\partial^2\psi}{\partial t^2}
},
\]
which is precisely the one-dimensional wave equation.

\subsection{(v) Wavefunction collapse and the double-slit experiment}

In an electron double-slit experiment, electrons are detected individually on a
screen, but after many electrons have passed through the apparatus, an
interference pattern develops. The electrons therefore exhibit wave-like
behaviour.

The de Broglie wavelength of the electrons determines the characteristic
diffraction scale. If the electrons are accelerated through a voltage $V$, their
classical kinetic energy is
\[
E=eV,
\]
and, using
\[
E=\frac{p^2}{2m},
\]
their wavelength is
\[
\lambda
=
\frac{h}{p}
=
\frac{h}{\sqrt{2meV}}.
\]

The double slit produces an interference pattern because the wavefunction
contains contributions associated with both possible paths.

If the electron is instead observed in a way that determines which slit it
passed through, the two-path superposition is destroyed. The wavefunction
collapses onto the corresponding measurement eigenstate, and the interference
pattern disappears.

Thus the act of measurement changes the state being measured: before the
which-path measurement, the electron is described by a superposition of the
possible paths; after the measurement, it has a definite measured path.

\subsection{(vi) Classical kinetic energy in terms of momentum}

Classically,
\[
E=\frac12mv^2
\]
and
\[
p=mv.
\]

Therefore
\[
v=\frac pm.
\]

Substituting,
\[
E
=
\frac12m\left(\frac pm\right)^2
=
\frac12m\frac{p^2}{m^2}.
\]

Hence
\[
\boxed{E=\frac{p^2}{2m}}.
\]

For comparison, the relativistic energy--momentum relation is
\[
E^2=p^2c^2+m^2c^4,
\]
or
\[
\boxed{E=\sqrt{p^2c^2+m^2c^4}},
\]
which reduces to $E=pc$ for a photon.

\subsection{(vii) Neutron diffraction}

The atomic lattice spacing is
\[
s=2.15\times10^{-10}\ \mathrm m.
\]
We interpret the stated $30^\circ$ angular spacing as the angle of the first
diffraction maximum from the central maximum, consistent with the supplied
solution. Thus
\[
\theta_1=30^\circ.
\]

For first-order Bragg diffraction,
\[
s\sin\theta_1=n\lambda,
\]
with $n=1$. Therefore
\[
\lambda=s\sin30^\circ=\frac{s}{2},
\]
so
\[
\boxed{\lambda=\frac{s}{2}}.
\]

Thus
\[
\lambda
=
\frac{2.15\times10^{-10}}{2}
=
1.08\times10^{-10}\ \mathrm m.
\]

Hence
\[
\boxed{\lambda\approx1.08\times10^{-10}\ \mathrm m}.
\]

Using
\[
p=\frac h\lambda
\]
and
\[
E=\frac{p^2}{2m_n},
\]
we obtain
\[
E
=
\frac{h^2}{2m_n\lambda^2}.
\]

Numerically,
\[
E\approx0.071\ \mathrm{eV}.
\]

Therefore
\[
\boxed{E\approx0.071\ \mathrm{eV}}.
\]

The neutron speed follows from
\[
p=m_nv,
\]
so
\[
v=\frac{h}{m_n\lambda}.
\]

Hence
\[
\boxed{v\approx3.68\times10^3\ \mathrm{m\,s^{-1}}}.
\]

The supplied solution also compares this energy with the thermal energy scale
\[
\frac32k_BT
\]
at room temperature and notes that the neutron energy is of the same order.

\subsection{(viii) Photoelectric effect and an equal de Broglie wavelength}

For silver,
\[
\phi=4.3\ \mathrm{eV},
\]
and the maximum photoelectron kinetic energy is
\[
K_{\max}=5.1\ \mathrm{eV}.
\]

The photoelectric equation is
\[
K_{\max}=\frac{hc}{\lambda}-\phi.
\]

Therefore
\[
\lambda
=
\frac{hc}{K_{\max}+\phi}.
\]

Using
\[
hc\approx1240\ \mathrm{eV\,nm},
\]
we find
\[
\lambda
=
\frac{1240}{5.1+4.3}
\approx131.9\ \mathrm{nm}.
\]

Thus
\[
\boxed{\lambda\approx131.9\ \mathrm{nm}}.
\]

Now consider an electron with the same de Broglie wavelength. Its kinetic energy is
\[
E_e=\frac{p^2}{2m_e},
\]
with
\[
p=\frac h\lambda.
\]

Hence
\[
E_e
=
\frac{h^2}{2m_e\lambda^2}.
\]

Substitution gives
\[
\boxed{E_e\approx8.65\times10^{-5}\ \mathrm{eV}}.
\]

Using the wavelength above gives
\[
E_e\approx8.65\times10^{-5}\ \mathrm{eV},
\]
consistent with the supplied handwritten solution. This illustrates the enormous
difference between the energy of a UV photon and the kinetic energy of a
non-relativistic electron having the same wavelength.

\subsection{(ix) Electron microscope resolution}

The resolution is approximately the electron wavelength:
\[
\lambda\approx\Delta x.
\]

Given
\[
\Delta x=0.42\ \mathrm{nm},
\]
we therefore require
\[
\lambda\approx0.42\ \mathrm{nm}.
\]

Using
\[
E=\frac{h^2}{2m_e\lambda^2},
\]
we obtain
\[
E
=
\frac{(6.626\times10^{-34})^2}
{2(9.109\times10^{-31})(0.42\times10^{-9})^2}.
\]

Converting to electron-volts,
\[
\boxed{E\approx8.53\ \mathrm{eV}}.
\]

Equivalently, if the electron is accelerated through a voltage $V$, then
\[
E=eV,
\]
and the non-relativistic relation becomes
\[
\boxed{
\frac{\lambda}{\mathrm{nm}}
\approx
\frac{1.23}{\sqrt{V}}
},
\]
where $V$ is measured in volts.

For larger accelerating voltages, relativistic corrections become important. The
supplied solution gives
\[
\boxed{
\frac{\lambda}{\mathrm{nm}}
\approx
\frac{1.23}
{\sqrt{V}\sqrt{1+9.78\times10^{-7}V}}
}
\]
as the corresponding relativistic expression.

\subsection{(x) Time and spatial extent of beta decay}

The emitted electron has energy
\[
\Delta E=5\ \mathrm{MeV}.
\]

From the energy-time uncertainty principle,
\[
\Delta E\,\Delta t\gtrsim\frac{\hbar}{2},
\]
so
\[
\Delta t
\gtrsim
\frac{\hbar}{2\Delta E}.
\]

Therefore
\[
\boxed{\Delta t\gtrsim6.58\times10^{-23}\ \mathrm s}.
\]

For the spatial extent, the problem assumes
\[
\Delta p\approx\sqrt{2m_e\Delta E}.
\]

Thus
\[
\Delta x
\gtrsim
\frac{\hbar}{2\Delta p}
=
\frac{\hbar}
{2\sqrt{2m_e\Delta E}}.
\]

Substitution gives
\[
\boxed{\Delta x\gtrsim4.36\times10^{-14}\ \mathrm m}.
\]

This is a nuclear-scale distance, consistent with the short-range nature of the
process.

% ============================================================
\section{Question 2 --- Particle in an infinite potential box}

A particle of mass $m$ is confined to
\[
0\le x\le a
\]
with
\[
V(x)=
\begin{cases}
0, & 0<x<a,\\
\infty, & x\le0\ \text{or}\ x\ge a.
\end{cases}
\]

\subsection{(i) Standing-wave solutions}

The infinite potential walls mean that the particle cannot exist outside the
interval $0<x<a$. Therefore the wavefunction must vanish at the walls:
\[
\psi(0,t)=0,
\qquad
\psi(a,t)=0.
\]

A standing wave satisfying the first boundary condition can be written as
\[
\psi(x,t)
=
A_n\sin\left(\frac{n\pi x}{a}\right)\cos\omega t.
\]

At $x=a$,
\[
\sin(n\pi)=0
\]
only when
\[
\boxed{n=1,2,3,\ldots}.
\]

Thus $n$ must be a positive integer.

The wavelength is related to the box length by
\[
\frac{n\lambda_n}{2}=a,
\]
so
\[
\lambda_n=\frac{2a}{n}.
\]

Therefore
\[
k_n=\frac{2\pi}{\lambda_n}
=
\frac{n\pi}{a}.
\]

The standing-wave picture is natural because the infinitely high walls impose
fixed nodes at the two ends of the box.

\subsection{(ii) Energy eigenvalues}

Inside the box,
\[
V=0,
\]
so the time-independent Schrödinger equation is
\[
-\frac{\hbar^2}{2m}\frac{\partial^2\psi}{\partial x^2}
=
E\psi.
\]

For
\[
\psi=A_n\sin\left(\frac{n\pi x}{a}\right),
\]
we have
\[
\frac{\partial^2\psi}{\partial x^2}
=
-\left(\frac{n\pi}{a}\right)^2\psi.
\]

Therefore
\[
\frac{\hbar^2}{2m}
\left(\frac{n\pi}{a}\right)^2\psi
=
E\psi.
\]

Hence
\[
\boxed{
E_n=\frac{\hbar^2\pi^2n^2}{2ma^2}
}.
\]

For a proton,
\[
m=m_p=1.6726\times10^{-27}\ \mathrm{kg},
\]
and
\[
a=0.47\times10^{-15}\ \mathrm m.
\]

Substitution gives
\[
E_n\approx927.2\,n^2\ \mathrm{MeV}.
\]

Thus for the ground state,
\[
\boxed{E_1\approx927.2\ \mathrm{MeV}}.
\]

The proton rest-mass energy is
\[
m_pc^2\approx938.4\ \mathrm{MeV}.
\]

Therefore
\[
\boxed{\frac{E_1}{m_pc^2}\approx0.988}.
\]

The ground-state confinement energy is therefore almost equal to the proton
rest-mass energy. The non-relativistic Schrödinger result is consequently outside
its regime of validity, but the numerical comparison is exactly what the
question asks for.

\subsection{(iii) Time dependence and normalisation}

The factor $\cos\omega t$ is a real representation of a standing-wave oscillation,
but the Schrödinger wavefunction is generally complex. For a stationary energy
eigenstate, the correct time dependence is
\[
e^{-iEt/\hbar}.
\]

Thus
\[
\psi_n(x,t)
=
A_ne^{-iE_nt/\hbar}
\sin\left(\frac{n\pi x}{a}\right).
\]

The Born interpretation requires
\[
\int_{-\infty}^{\infty}|\psi|^2\,dx=1.
\]

Since the wavefunction is zero outside the box,
\[
\int_0^a|\psi|^2\,dx=1.
\]

The phase factors cancel:
\[
|\psi|^2
=
A_n^2
\sin^2\left(\frac{n\pi x}{a}\right).
\]

Hence
\[
A_n^2
\int_0^a
\sin^2\left(\frac{n\pi x}{a}\right)\,dx
=1.
\]

Using
\[
\int_0^a
\sin^2\left(\frac{n\pi x}{a}\right)\,dx
=
\frac a2,
\]
we find
\[
A_n^2\frac a2=1,
\]
so
\[
\boxed{A_n=\sqrt{\frac2a}}.
\]

Therefore the normalised stationary state is
\[
\boxed{
\psi_n(x,t)
=
\sqrt{\frac2a}\,
e^{-iE_nt/\hbar}
\sin\left(\frac{n\pi x}{a}\right)
}
\]
for $0<x<a$.

\subsection{(iv) Momentum expectation values}

A stationary standing wave is a superposition of waves travelling in opposite
directions. Consequently the average momentum is zero:
\[
\boxed{\langle p\rangle=0}.
\]

The energy is
\[
E=\frac{p^2}{2m},
\]
so
\[
p^2=2mE.
\]

Using
\[
E_n=\frac{\hbar^2\pi^2n^2}{2ma^2},
\]
we obtain
\[
\langle p^2\rangle
=
2mE_n
=
\frac{\hbar^2\pi^2n^2}{a^2}.
\]

Thus
\[
\boxed{
\langle p^2\rangle
=
\frac{\hbar^2\pi^2n^2}{a^2}
}.
\]

Since $\langle p\rangle=0$,
\[
\Delta p
=
\sqrt{\langle p^2\rangle-\langle p\rangle^2}
=
\frac{\hbar n\pi}{a}.
\]

\subsection{(v) Position expectation values}

The expectation value of position is
\[
\langle x\rangle
=
\int_0^a x|\psi|^2\,dx.
\]

Using the normalised wavefunction,
\[
\langle x\rangle
=
\frac2a
\int_0^a
x\sin^2\left(\frac{n\pi x}{a}\right)\,dx.
\]

Using
\[
\int x\sin^2(\alpha x)\,dx
=
\frac{x^2}{4}
-\frac{x\sin(2\alpha x)}{4\alpha}
-\frac{\cos(2\alpha x)}{8\alpha^2}
+C,
\]
with
\[
\alpha=\frac{n\pi}{a},
\]
and the identities
\[
\sin(2n\pi)=0,
\qquad
\cos(2n\pi)=1,
\]
we obtain
\[
\boxed{\langle x\rangle=\frac a2}.
\]

For the second moment,
\[
\langle x^2\rangle
=
\frac2a
\int_0^a
x^2\sin^2\left(\frac{n\pi x}{a}\right)\,dx.
\]

Using the corresponding integral supplied on the problem sheet gives
\[
\boxed{
\langle x^2\rangle
=
\frac{a^2}{3}
\left(
1-\frac{3}{2n^2\pi^2}
\right)
}.
\]

Therefore
\[
(\Delta x)^2
=
\langle x^2\rangle-\langle x\rangle^2.
\]

Hence
\[
(\Delta x)^2
=
\frac{a^2}{3}
\left(
1-\frac{3}{2n^2\pi^2}
\right)
-\frac{a^2}{4}.
\]

Simplifying,
\[
(\Delta x)^2
=
\frac{a^2}{12}
-\frac{a^2}{2n^2\pi^2}.
\]

Thus
\[
\boxed{
\Delta x
=
\frac{a}{2\pi}
\sqrt{\frac{n^2\pi^2}{3}-2}
}.
\]

\subsection{(vi) Verification of the uncertainty principle}

From part (iv),
\[
\Delta p
=
\frac{\hbar n\pi}{a}.
\]

Combining this with
\[
\Delta x
=
\frac{a}{2\pi}
\sqrt{\frac{n^2\pi^2}{3}-2},
\]
we obtain
\[
\Delta x\,\Delta p
=
\frac{\hbar}{2}
\sqrt{\frac{n^2\pi^2}{3}-2}.
\]

Therefore
\[
\boxed{
\Delta x\,\Delta p
=
\frac{\hbar}{2}
\sqrt{\frac{n^2\pi^2}{3}-2}
}.
\]

The smallest possible value occurs for the smallest allowed quantum number,
\[
n=1.
\]

Then
\[
\sqrt{\frac{\pi^2}{3}-2}
\approx1.14.
\]

Hence
\[
\Delta x\,\Delta p
\approx1.14\frac{\hbar}{2}
>
\frac{\hbar}{2}.
\]

Therefore every allowed particle-in-a-box state satisfies
\[
\boxed{
\Delta x\,\Delta p\ge\frac{\hbar}{2}
},
\]
which verifies the Heisenberg uncertainty principle.

% ============================================================
\section{Question 3 --- Entangled photons and quantum cryptography}

Two vertically polarised entangled photons travel in opposite directions to
detectors $A$ and $B$. The detectors are oriented at angles $\theta$ and $\phi$.

The problem asks first for a classical calculation in which the measurements are
treated as independent, and then for the Copenhagen interpretation in which the
first measurement affects the state of the second photon.

\subsection{Classical calculation}

Using Malus' law, the probabilities for detector $A$ are
\[
P(X_A)=\cos^2\theta,
\qquad
P(Y_A)=\sin^2\theta.
\]

Similarly,
\[
P(X_B)=\cos^2\phi,
\qquad
P(Y_B)=\sin^2\phi.
\]

The two detectors give different outcomes in the two cases
\[
X_A,Y_B
\]
and
\[
Y_A,X_B.
\]

Assuming the measurements are independent,
\[
P_{\rm classical}(\mathrm{mismatch})
=
P(X_A)P(Y_B)+P(Y_A)P(X_B).
\]

Therefore
\[
\boxed{
P_{\rm classical}
=
\cos^2\theta\sin^2\phi
+
\sin^2\theta\cos^2\phi
}.
\]

For
\[
\theta=-45^\circ,
\qquad
\phi=30^\circ,
\]
we have
\[
\cos^2(-45^\circ)=\sin^2(-45^\circ)=\frac12,
\]
and
\[
\sin^2(30^\circ)=\frac14,
\qquad
\cos^2(30^\circ)=\frac34.
\]

Hence
\[
P_{\rm classical}
=
\frac12\frac14
+
\frac12\frac34
=
\frac12.
\]

Thus
\[
\boxed{P_{\rm classical}(\mathrm{mismatch})=\frac12}.
\]

\subsection{Copenhagen interpretation}

Suppose detector $A$ measures first.

If $A$ measures $X_A$, the entangled state collapses so that the corresponding
polarisation of the photon arriving at $B$ is referenced to the $A$ measurement
axis. The angle between the $A$ and $B$ measurement axes is
\[
\phi-\theta.
\]

Thus the conditional probabilities for detector $B$ involve
\[
\cos^2(\phi-\theta)
\]
and
\[
\sin^2(\phi-\theta).
\]

There are two possibilities for the first measurement:
\[
P(X_A)=\cos^2\theta,
\qquad
P(Y_A)=\sin^2\theta.
\]

The probability of a mismatch is therefore
\[
P_{\rm quantum}(\mathrm{mismatch})
=
\cos^2\theta\sin^2(\phi-\theta)
+
\sin^2\theta\cos^2(\phi-\theta).
\]

This simplifies to
\[
\boxed{
P_{\rm quantum}(\mathrm{mismatch})
=
\sin^2(\phi-\theta)
}.
\]

For
\[
\theta=-45^\circ,
\qquad
\phi=30^\circ,
\]
the relative angle is
\[
\phi-\theta
=
30^\circ-(-45^\circ)
=
75^\circ.
\]

Therefore
\[
P_{\rm quantum}
=
\sin^2(75^\circ).
\]

Using
\[
\sin75^\circ
=
\sin(30^\circ+45^\circ)
=
\frac{1}{2\sqrt2}(1+\sqrt3),
\]
we obtain
\[
P_{\rm quantum}
=
\frac{(1+\sqrt3)^2}{8}
\approx0.933.
\]

Hence
\[
\boxed{
P_{\rm quantum}(\mathrm{mismatch})\approx0.933
}.
\]

The ratio of the quantum prediction to the classical prediction is
\[
\frac{P_{\rm quantum}}{P_{\rm classical}}
=
\frac{(1+\sqrt3)^2/8}{1/2}
=
\frac{(1+\sqrt3)^2}{4}
\approx1.87.
\]

Thus
\[
\boxed{
\frac{P_{\rm quantum}}{P_{\rm classical}}\approx1.87
}.
\]

The quantum result is therefore substantially different from the independent
classical prediction. This difference is the basis of the quantum-cryptographic
idea described in the problem sheet: the statistics of polarisation measurements
can reveal whether an eavesdropper has disturbed the transmitted quantum states.

\end{document}
